\begin{proof}
The even-shift translation and its quadratic correction are those of \citet[Proposition~7.8, pp.~91--92, and Proposition~C.2, p.~102]{DuEtAl2026v1}. Their representatives in the fixed root and cusp charts are as follows.
The global roots exist because the zero orders of \(t\) and \(t-1\)
on the simply connected uniformizing surface are divisible by three
and four. The local characters and their coprime orders give
\[
 \ell(g_1z)-\ell(z)=-\frac13,\qquad
 \ell(g_2z)-\ell(z)=\frac14,\qquad
 \ell(g_0z)-\ell(z)=\frac1{12}\pmod{\mathbb Z}.
\]
A change of logarithm adds an integral \(h\widehat\delta\) period,
so \eqref{mapsI:closing:eq:C-explicit-closing} is a well-defined torus translation.
Since \(\Pi_{c+h}=\Pi_c(I+hE)\), the two ordinary lattices agree;
the source-to-target lattice map of the identity linear part is
\(I-hE\).

For compatible logarithm lifts, covariance gives the affine defects
\[
 b_{h,c}(g_jz)-R_j(z)b_{h,c}(z)
       =\Pi_c(g_jz)\eta_j,\qquad
 \eta_1=\frac h2(\widehat u-\widehat w-\widehat\delta),\quad
 \eta_2=\frac h2(-\widehat w+\widehat\delta).
\]
Their \(\gamma\)-coordinates vanish, so the same right sides use
\(\Pi_{c+h}\). The vectors are integral for even \(h\). At the
cusp the defect is zero: \((A_0-I)\widehat\gamma=-\widehat\delta\)
cancels the increment of \(\ell\). These equations prove ordinary
and cusp descent with the stated orientation of the peripheral
loops.

For the finite comparison put \(w_j=(2\pi i)^{-1}\log s_j\).
The functions \(F/(s_1G)\) and \(Fs_2/G\) are invariant
holomorphic units on the respective full root discs. Choose
holomorphic logarithms \(\kappa_j\) of them. Then
\[
 \ell=\epsilon_jw_j+\kappa_j\pmod{\mathbb Z},\qquad
 \epsilon_1=1,\quad\epsilon_2=-1.
\]
The exact finite-to-ordinary comparison satisfies
\(\sigma_{j,c+h}-\sigma_{j,c}=h\epsilon_jw_je\).
Consequently the required finite-chart translation is
\[
 b^{(j)}_{h,c}
 =-\frac h{12}\Pi_c\widehat\gamma
       +\left(h\kappa_j-\frac{h^2}{24}\right)e.
\]
It is holomorphic on the full root disc and satisfies
\[
 b^{(j)}_{h,c}(g_jz)-R_j(z)b^{(j)}_{h,c}(z)
       -\frac{h\epsilon_j}{m_j}e
 =\Pi_c(g_jz)\eta_j.
\]
This is precisely equivariance for the source and target affine
finite actions, including the changed period column. The exact
identity
\(\sigma_{j,c+h}+b^{(j)}_{h,c}
=\sigma_{j,c}+b_{h,c}\) proves overlap agreement. 

At the cusp let \(q=e^{2\pi i s}\) and
\(C_{c,21}=\beta^{\rm part}+c+s\), which is holomorphic at
\(q=0\). The function
\[
 \widehat H=e^{2\pi i s/12}F/G
 \quad\hbox{satisfies}\quad
 \widehat H^{12}=q\frac{t^4}{(t-1)^3}.
\]
The right side is a unit because \(t\) has a simple pole at the
cusp. The character formula makes \(\widehat H\) single-valued;
it therefore extends as a nonvanishing holomorphic function. Write
\(\ell=-s/12+\kappa_0(q)\), with \(\kappa_0\) holomorphic.
Since
\(\Pi_c\widehat\gamma=(6\mu,-s+C_{c,21})^t\),
the logarithmic terms cancel in \eqref{mapsI:closing:eq:C-explicit-closing}:
\[
 \bar b_{h,c}(q)=
 \left(-\frac h2\mu(q),\,
 -\frac h{12}C_{c,21}(q)+h\kappa_0(q)-\frac{h^2}{24}\right)^t.
\]
Exponentiating gives a holomorphic unit in \((\mathbb C^*)^2\).
On every toric chart its action multiplies a character coordinate
by an integral character of this unit, so it and its inverse
extend across every double curve and both triple points. The
source and target cusp lattice multipliers differ by
\((1,e^{2\pi i h\lambda_1})=(1,1)\), for
\(\lambda\in\mathbb Z^2\). They are therefore identical. The
unit multiplication commutes with their toric shears and descends.
Exponentiation of the displayed formula gives exactly the ordinary
translation on the punctured cusp. We have now verified all
three comparisons and their inverses.

The uncorrected translations have composition defect
\(-hk e/12\). The constants in
\eqref{mapsI:closing:eq:C-explicit-closing} cancel it, since
\((h+k)^2-h^2-k^2=2hk\). Hence
\[
 \widetilde F_{k,c+h}\widetilde F_{h,c}
       =\widetilde F_{h+k,c},\qquad
 \widetilde F_{0,c}=\mathrm{id}.
\]
Composing this even-shift map with the inverse of the fixed choice in Proposition~\ref{fam:src:A.3.1} gives an automorphism over the identity of the base. The full automorphism theorem, Theorem~\ref{night:P5.5}, identifies it with a scalar in $\mathbb C^*$. A path from that scalar to $1$ supplies an isotopy, so the two closing choices define the same smooth monodromy class.

For \(h=12\), remove the period
\(-\Pi_c\widehat\gamma\) from \eqref{mapsI:closing:eq:C-explicit-closing}.
The remaining translation is \((12\ell-6)e\). Its exponential
is \((F/G)^{12}=K(t)\), proving
\eqref{mapsI:closing:eq:C-twelve-scalar}. The finite corrections divide it by
\(s_1^{12}\) and multiply it by \(s_2^{12}\). With the original
uniformizers, these are units because
\(t/s_1^3\) and \((t-1)/s_2^4\) are units. 

Put \(H(q)=qK(t(q))=e^{24\pi i\kappa_0(q)}\). The cusp unit
for \(h=12\) is exactly
\[
 U_{12,c}(q)=
 \left(e^{-12\pi i\mu(q)},\,
       e^{-2\pi i C_{c,21}(q)}H(q)\right).
\]
The \(\widehat\gamma\) deck transformation has multiplier
\[
 \left(e^{12\pi i\mu(q)},\,
       e^{2\pi i C_{c,21}(q)}q^{-1}\right),
\]
because \(B_0(1,0)^t=(0,-1)^t\). Multiplying these two
expressions gives \((1,q^{-1}H)=(1,K)\). This proves the
claimed comparison, including its sign, on the entire cusp
quotient. The extending map is represented by $U_{12,c}$, and the displayed deck multiplier identifies it with the scalar representative on the punctured chart.

Finally \(w=e^{-\pi i c}\) makes six positive loops the path
\(c(r)=c_0-12r\); its endpoint is closed by the positive shift
\(h=12\). For parallel transport $P$ along that path the closed diffeomorphism is therefore $\widetilde F_{12,c_0-12}\circ P$. The overlap difference between the scalar logarithm and the two-coordinate unit logarithm is exactly $\Pi_c\widehat\gamma$. Scaling both logarithms by a real number $a$ changes this difference to $a\Pi_c\widehat\gamma$. For $0<a<1$, this vector is not a period: the real period map is an isomorphism and $a\widehat\gamma\notin\Lambda$. Thus the separately scaled chart maps fail the overlap equation at precisely this fractional period.
\end{proof}